\begin{agradecimentos}
À minha esposa, Jéssica Alves de Souza, pelo apoio incondicional, pela paciência e por ter sido o principal incentivo para que eu retornasse aos estudos na Universidade Tecnológica Federal do Paraná. Sem o seu amor, companheirismo e encorajamento constante, a realização deste trabalho e esta conquista não teriam sido possíveis.

À minha orientadora, Prof.ª Dra. Sediane Carmem Lunardi Hernandes, por acreditar no potencial deste trabalho e por insistir para que uma ideia inicial de simulador de comandos se transformasse em um Trabalho de Conclusão de Curso estruturado e de impacto real para a formação acadêmica.

Ao Prof. Hermano Pereira, pelos ensinamentos aprofundados no ecossistema Linux aplicados às disciplinas de redes de computadores, cuja exigência técnica e rigor prático elevaram o meu domínio da linha de comando e da infraestrutura de sistemas.

Ao Prof. Luciano Ogiboski, pelo suporte e pela orientação fundamentais que viabilizaram o meu retorno à UTFPR.

Ao Prof. Emerson, por compartilhar sua visão empreendedora e empresarial, me incentivando nos primeiros passos corporativos e abrindo portas para a contratação dos primeiros estagiários da minha empresa.

Aos demais professores do Departamento Acadêmico de Informática (DAINF) e colaboradores da UTFPR Câmpus Guarapuava, que contribuíram de forma direta ou indireta para a minha trajetória acadêmica e profissional ao longo deste curso.
\end{agradecimentos}
